\clearpage
\begingroup
\HistoricalHead{249}{R}{BEMERKUNGEN.}
\HistoricalFoot{XI1.}{32}
\makeatletter\tagsleft@true\let\veqno\leqno\makeatother
\begin{center}\textbf{BEMERKUNGEN.}\end{center}
\noindent\hspace*{\parindent}Die ältesten im Nachlaß erhaltenen \textsc{Gauss}schen Rechnungen und Notizen zur Bahnbestimmung, insbesondere für Ceres, finden sich in den Schedae Ag und Ah, sowie im Handbuch Bb, und zwar in den ersten beiden vorzugsweise die numerischen Rechnungen, im letzteren dagegen nur Formelentwicklungen; sie stammen aus dem November 1801, sodaß also die allerersten Untersuchungen aus dem September und Oktober (vergl. das \textit{Tagebuch} No. 119 u. 121, Werke X1, S. 561, 563) verloren gegangen sind. Man wird aber in der Notiz [I.], oben S. 221ff., die auf den ersten Seiten des Handbuchs Bb unter der Überschrift »Erste Abhandlung« steht, einen Niederschlag der in der Tagebuchnotiz No. 121 erwähnten »permultae formulae« erblicken dürfen.

Die Notiz [II.], oben S. 232ff., steht in der Scheda Ag. Die Zahlenbeispiele bilden die fortlaufende Rechnung zur Ableitung der I. Elemente der Ceres, die \textsc{Gauss} an v. \textsc{Zach} schickte (Werke VI, S. 200f.). Die vorausgesetzten Werte von $i$ und $\AstNode$ sind vielleicht einer noch früheren, verloren gegangenen Originalrechnung\textsuperscript{***)}\begingroup\renewcommand{\thefootnote}{***)}\footnotetext{\textsc{Gauss} sagt in der Einleitung zur \textit{Theoria motus} (Werke VII, 1906, S. 8), daß er die erste Bahnbestimmung der Ceres im Oktober 1801 gemacht habe.}\endgroup\ entnommen; bemerkenswert ist, daß \textsc{Gauss} von der hier angewandten »Dritten Verbesserungsmethode« (siehe S. 229) später in der \textit{Theoria motus}, art. 82 (vergl. auch art. 126) sagt, daß sie »orbitae dimensiones eruendi magnam praecisionem nunquam admittet, nisi tria loca heliocentrica intervallis considerabilibus ab invicem distent«.
\clearpage
\HistoricalHead{250}{L}{THEORETISCHE ASTRONOMIE. NACHLASS.}
\SourceLeaf{255}
Die Notiz [III.], oben S. 241ff., steht in der Scheda Ah; es ist dies das Original der Berechnung der Ceres-Elemente III. durch Verbesserung der Elemente II. Die numerischen Rechnungen sind beim Abdruck gekürzt worden, insbesondere wurden die bei drei Beobachtungen gleichlautenden Rechnungen nur bei der ersten Beobachtung ausführlich abgedruckt.

Wie \textsc{Gauss} nicht nur in der Einleitung zur \textit{Theoria motus} und in einem Briefe an \textsc{Olbers} vom 6. August 1802 sagt, sondern wie auch aus der Anmerkung hervorgeht, mit der \textsc{von Lindenau} die Herausgabe des ihm von \textsc{Gauss} übergebenen Manuskripts \textit{Summarische Übersicht der zur Bestimmung der Bahnen der beiden neuen Hauptplaneten angewandten Methoden} (Monatl. Corresp., Sept. 1809 und Werke, Bd. VI, S. 148) begleitet hat, unterscheiden sich die Methoden der \textit{Theoria motus} sehr wesentlich von den ältesten \textsc{Gauss}schen Rechnungsweisen; die ersteren sind nach den Tagebuchnotizen No. 125—129 (Werke X1, S. 564, 565) auch erst in den Jahren 1805/06 entstanden. Doch erwähnt \textsc{Gauss} auch einige der letzteren gelegentlich in der \textit{Theoria motus}; auf die betreffenden Parallelstellen ist oben durch Fußnoten hingewiesen. Die erwähnte \textit{Summarische Übersicht} hatte \textsc{Gauss} schon im Jahre 1802 zusammengestellt und am 6. August an \textsc{Olbers} gesandt (siehe \textsc{Wilhelm Olbers}, Sein Leben und seine Werke II, 1, 1900, S. 65).

In den Artikeln 124—129 der \textit{Theoria motus} bespricht \textsc{Gauss} zehn verschiedene Methoden der Bahnbestimmung; man wird bemerken, daß die unter [I.] abgedruckten Notizen [3]—[6] sich an die erste dieser Methoden, die Notiz [7] an die siebente (\textit{Theoria motus}, art. 129, No. II) und die Notiz [8] an die achte (\textit{Theoria motus}, art. 129, No. III) anschließt; doch fehlt bei der Notiz [8] der Hinweis, daß die Interpolation auf die verbesserten Werte durch Vergleichung der berechneten Zwischenzeiten mit den beobachteten erfolgen soll. Die Notiz [II.] schließt sich ebenfalls an die achte Methode und die Notiz [III.] an die siebente an.

Über die Ableitung der Formeln findet man das nötige teils in der \textit{Summarischen Übersicht} teils in der \textit{Theoria motus} an den in den Fußnoten angegebenen Stellen; nur für die Formeln a—d auf S. 222—223, die streng gelten, scheint es nicht überflüssig hier eine kurze Ableitung zu geben:

Sind $x,y,z$ die heliozentrischen Koordinaten des Planeten, $X,Y,Z$ die Koordinaten der Sonne in Bezug auf die Erde, $U$, wie bei \textsc{Gauss}, die mittlere Länge der Sonne und $A$ die halbe große Achse der Erdbahn, so hat man:
\begin{equation*}
\frac{d^2x}{dU^2}+\frac{A^3x}{r^3}=0,\qquad
\frac{d^2X}{dU^2}+\frac{A^3X}{R^3}=0;\tag*{(1)}
\end{equation*}
\noindent da aber
\begin{equation*}
x=\Delta\cos\alpha-R\cos V,\quad X=R\cos V,\quad x+X=\Delta\cos\alpha,\tag*{(2)}
\end{equation*}
\noindent so ergibt die Summe der Gleichungen (1)
\begin{equation*}
\begin{aligned}
\frac{d^2(\Delta\cos\alpha)}{dU^2}&=-A^3\left(\frac{x}{r^3}+\frac{X}{R^3}\right)\\
&=-\frac{A^3\Delta\cos\alpha}{r^3}-A^3R\cos V\left(\frac1{R^3}-\frac1{r^3}\right)
\end{aligned}\tag*{(3)}
\end{equation*}
\noindent und ebenso erhält man aus den den Gleichungen (1) entsprechenden Gleichungen in $y,Y$ und $z,Z$, sowie aus den den Gleichungen (2) entsprechenden:
\[
\begin{array}{lll}
y=\Delta\sin\alpha-R\sin V,&Y=R\sin V,&y+Y=\Delta\sin\alpha,\\
z=\Delta\mathbin{.}\theta,&Z=0,
\end{array}
\]
\noindent die folgenden Parallelgleichungen zu (3):
\begin{equation*}
\frac{d^2(\Delta\sin\alpha)}{dU^2}=-\frac{A^3\Delta\sin\alpha}{r^3}-A^3R\sin V\left(\frac1{R^3}-\frac1{r^3}\right),\tag*{(4)}
\end{equation*}
\begin{equation*}
\frac{d^2(\Delta\mathbin{.}\theta)}{dU^2}=-\frac{A^3\Delta\mathbin{.}\theta}{r^3}.\tag*{(5)}
\end{equation*}
\clearpage
\HistoricalHead{251}{R}{BEMERKUNGEN.}
\SourceLeaf{256}
Führt man die Differentiation auf den linken Seiten der Gleichungen (3)—(5) aus, multipliziert sie zur Elimination von $\dfrac{d^2\Delta}{dU^2}$ und $\dfrac{d\Delta}{dU}$ der Reihe nach mit
\[
\theta\frac{d\alpha}{dU}\cos\alpha-\frac{d\theta}{dU}\sin\alpha,\quad
\theta\frac{d\alpha}{dU}\sin\alpha+\frac{d\theta}{dU}\cos\alpha,\quad
-\frac{d\alpha}{dU}
\]
\noindent und addiert, so wird
\begin{equation*}
\begin{aligned}
&\Delta\frac{d^2\alpha}{dU^2}\frac{d\theta}{dU}-\Delta\frac{d\alpha}{dU}\frac{d^2\theta}{dU^2}-\Delta\mathbin{.}\theta\left(\frac{d\alpha}{dU}\right)^3\\
&\qquad=-A^3R\left(\frac1{R^3}-\frac1{r^3}\right)\left(\theta\frac{d\alpha}{dU}\cos(V-\alpha)+\frac{d\theta}{dU}\sin(V-\alpha)\right)
\end{aligned}\tag*{(6)}
\end{equation*}
\noindent und dies ist die Gleichung (a) der Nachlaßnotiz, wo nur $A=1$ gesetzt ist. Man vergleiche hierzu den Brief von \textsc{Olbers} an \textsc{Gauss} vom 11. September 1802 und \textsc{Laplace}, \textit{Mécanique céleste}, Première Partie, Livre II, No. 31 (Band I, S. 207 der Originalausgabe von 1799).

Um die Gleichung (b) abzuleiten, multipliziert man in ähnlicher Weise die Gleichungen (3)—(5) zur Elimination von $\dfrac{d^2\Delta}{dU^2}$ und $R$ der Reihe nach mit $-\theta\sin V$, $\theta\cos V$, $\sin(V-\alpha)$ und addiert.

Die Gleichungen (a) und (b) lassen sich auch aus I und II durch den Grenzübergang $\lim\Delta U=0$ ableiten, wobei noch zu beweisen ist, daß die in I und II vernachlässigten Glieder beim Grenzübergang verschwinden.

Die Gleichung (c) erhält man durch Multiplikation der Gleichung (3) bezw. (4) mit $-\sin\alpha$ bezw. $\cos\alpha$ und Addition und endlich die Gleichung (d) durch Multiplikation der Gleichungen (3)—(5) mit $-\theta\cos\alpha$, $-\theta\sin\alpha$, $+1$ und Addition.

Die Gleichungen für $\log\Delta$ und $\log\Delta''$, S. 223, beruhen auf der Voraussetzung, daß die Logarithmen der Abstände von der Erde sich proportional den Zwischenzeiten ändern, daß also
\[
\frac{\log\Delta''-\log\Delta'}{t''-t'}=\frac{\log\Delta'-\log\Delta}{t'-t}.
\]
\noindent\hspace*{\parindent}Mit der auf S. 229 bei Gelegenheit der »zweiten Verbesserungsmethode« erwähnten »Aufgabe« hat sich \textsc{Gauss} öfters beschäftigt; es finden sich darüber Notizen im Handbuch Bb, S. 15 und in der Scheda Ah, S. 5 (mit Beispielen) und \textsc{Gauss} hat diesen Teil seiner Untersuchungen in der Monatl. Corresp. Juni 1802 (Bd. VI, S. 87) veröffentlicht; vergl. auch \textit{Theoria motus}, art. 74.

Vielleicht ist es erwünscht, die Ableitung der Formel für $p$ in der Fußnote S. 225 aus der vorhergehenden Gleichung hier zu geben, da sie nicht sofort zu übersehen ist. Die Größen $\delta$ und $\varphi$ sind klein, ihr Cosinus nahe gleich Eins; setzt man also
\[
\frac{\cos\frac12\delta}{\cos\varphi\left(1-\frac{2\sin\frac14\delta^2}{p\cos\varphi}\surd(rr'\cos2\varphi)\right)}=1+\xi,
\qquad\cos2\varphi=1+x,
\]
\noindent so hat man unter Vernachlässigung der Quadrate von $x$ und $\xi$ näherungsweise:
\[
\begin{aligned}
m\surd p&=\frac{\delta rr'}3(1-x+2(1+x)(1+2\xi))\\
&=\delta rr'(1+\tfrac13x+\tfrac43\xi)\\
&=\delta rr'\mathord{{}^{3}\!\surd}(1+x+4\xi)\\
&=\delta rr'\mathord{{}^{3}\!\surd}(\cos2\varphi\mathbin{.}(1+\xi)^4).
\end{aligned}
\]
\noindent\hspace*{\parindent}Der auf S. 230 (letzte Zeile) von \textsc{Gauss} angegebene Wert für die größte Mittelpunktsgleichung ist eine Näherungsformel ähnlicher Art, wie die vorstehende; \textsc{Gauss} bediente sich häufig solcher Näherungsausdrücke: man vergleiche Werke VII, 1906, S. 330, 332 u. a., im besonderen S. 335. Entwickelt man nach Potenzen von $e$, so ergibt die Näherungsformel die Reihe:
\[
2e+\frac{11}{48}e^3+\frac{528}{5120}e^5+\mathinner{\cdotp\cdotp\cdotp},
\]
\begin{flushright}32\textsuperscript{*}\end{flushright}
\clearpage
\HistoricalHead{252}{L}{THEORETISCHE ASTRONOMIE. NACHLASS.}
\SourceLeaf{257}
\noindent während der genaue Wert ist:
\[
2e+\frac{11}{48}e^3+\frac{599}{5120}e^5+\mathinner{\cdotp\cdotp\cdotp}.
\]
\noindent\hspace*{\parindent}Außer den oben beim Text angegebenen wurden noch einige weitere kleinere Schreibfehler beim Abdruck verbessert.

Es schien angebracht, im vorstehenden die ältesten Urkunden der Arbeiten von \textsc{Gauss} auf diesem Gebiete zum Abdruck zu bringen, wenn auch viele Einzelheiten aus der \textit{Summarischen Übersicht} und der \textit{Theoria motus} entnommen werden können.
\begin{flushright}\textsc{Brendel}.\end{flushright}
\begin{center}\rule{.24\linewidth}{.3pt}\end{center}
\clearpage
\endgroup
